\section{Билинейные и полуторалинейные формы. Теорема о представлении б.ф. в виде суммы симметричной и кососимметричной б.ф. Замечание о симметричных б.ф.}

\begin{definition}
Пусть $V$ --- линейное пространство над полем $F$ ($F = \mathbb{R}$ или $\mathbb{C}$). Функция $f: V \times V \to F$ называется \textbf{билинейной формой}, если она линейна по каждому из своих аргументов, то есть для любых $x, x_1, x_2, y, y_1, y_2 \in V$ и $\lambda, \mu \in F$ выполняются условия:
\begin{enumerate}
    \item $f(x_1 + x_2, y) = f(x_1, y) + f(x_2, y)$;
    \item $f(\lambda x, y) = \lambda f(x, y)$;
    \item $f(x, y_1 + y_2) = f(x, y_1) + f(x, y_2)$;
    \item $f(x, \mu y) = \mu f(x, y)$.
\end{enumerate}
\end{definition}

\begin{definition}
Билинейная форма $f$ называется:
\begin{itemize}
    \item \textbf{симметричной}, если $f(x, y) = f(y, x)$ для всех $x, y \in V$;
    \item \textbf{кососимметричной}, если $f(x, y) = -f(y, x)$ для всех $x, y \in V$.
\end{itemize}
\end{definition}

\begin{theorem}[О разложении билинейной формы]
Любая билинейная форма $f$ на линейном пространстве $V$ над полем $F$ (где $\operatorname{char} F \neq 2$) единственным образом представляется в виде суммы симметричной билинейной формы $f_1$ и кососимметричной билинейной формы $f_2$:
\[
f(x, y) = f_1(x, y) + f_2(x, y).
\]
\end{theorem}

\begin{proof}
\begin{enumerate}
    \item \textbf{Существование:} Определим формы $f_1$ и $f_2$ следующими равенствами:
    \[
    f_1(x, y) = \frac{1}{2} \big( f(x, y) + f(y, x) \big), \quad f_2(x, y) = \frac{1}{2} \big( f(x, y) - f(y, x) \big).
    \]
    Очевидно, что $f_1(x, y) + f_2(x, y) = f(x, y)$. Проверим свойства:
    \[
    f_1(y, x) = \frac{1}{2} \big( f(y, x) + f(x, y) \big) = f_1(x, y) \quad \text{(симметрична)},
    \]
    \[
    f_2(y, x) = \frac{1}{2} \big( f(y, x) - f(x, y) \big) = -f_2(x, y) \quad \text{(кососимметрична)}.
    \]
    Линейность $f_1$ и $f_2$ по обоим аргументам следует из линейности исходной формы $f$.
    
    \item \textbf{Единственность:} Пусть $f(x, y) = s(x, y) + k(x, y)$, где $s$ --- симметричная, а $k$ --- кососимметричная формы. Тогда:
    \[
    f(y, x) = s(y, x) + k(y, x) = s(x, y) - k(x, y).
    \]
    Складывая и вычитая эти два равенства, получаем:
    \[
    f(x, y) + f(y, x) = 2s(x, y) \implies s(x, y) = \frac{1}{2}\big(f(x, y) + f(y, x)\big) = f_1(x, y),
    \]
    \[
    f(x, y) - f(y, x) = 2k(x, y) \implies k(x, y) = \frac{1}{2}\big(f(x, y) - f(y, x)\big) = f_2(x, y).
    \]
    Разложение единственно.
\end{enumerate}
\hfill$\blacksquare$
\end{proof}

\begin{remark}[О симметричных билинейных формах]
Симметричная билинейная форма $f$ полностью определяется своими значениями на «диагонали», то есть значениями $f(x, x)$. Это следует из \textbf{тождества поляризации}:
\[
f(x+y, x+y) = f(x, x) + f(x, y) + f(y, x) + f(y, y).
\]
Поскольку $f$ симметрична ($f(x, y) = f(y, x)$), получаем:
\[
f(x+y, x+y) = f(x, x) + 2f(x, y) + f(y, y).
\]
Отсюда выражается $f(x, y)$:
\[
f(x, y) = \frac{1}{2} \big( f(x+y, x+y) - f(x, x) - f(y, y) \big).
\]
\end{remark}

\begin{definition}
Если $V$ --- линейное пространство над полем $\mathbb{C}$, функция $f: V \times V \to \mathbb{C}$ называется \textbf{полуторалинейной} (или эрмитовой) формой, если она линейна по первому аргументу и антилинейна по второму:
